\section{Evaluation and dependence of the second coefficient}
\label{sec:coefficient}

\subsection{The regularized Riemann sum}

\begin{proposition}[The exact critical sequence]
\label{prop:exact-law}
For each fixed $\kappa>0$, the integral in
\eqref{eq:coefficient-definition} converges and
\begin{equation}
 \sum_{n\ge1}J_\kappa(k/n)
           =k\log k+C(\kappa)k+o(k).
 \label{eq:ids-sum-law}
\end{equation}
Consequently \eqref{eq:main-law} holds for $b_n=\kappa n$.
\end{proposition}

\begin{proof}
For $x>0$, define
\begin{equation}
 G_\kappa(x)=J_\kappa(1/x)-x^{-1}\ind_{(0,1]}(x).
 \label{eq:regularized-summand}
\end{equation}
By \eqref{eq:cell-bounds}, $\abs{G_\kappa(x)}\le1$
for $0<x\le1$. By \eqref{eq:ids-gap}, it vanishes for
$x>1/r(\kappa)$. Lemma~\ref{lem:phase} shows that it is
continuous except possibly at the cutoff $x=1$, away from
the endpoint zero. It is thus absolutely integrable.

Its normalized Riemann sums obey
\begin{equation}
 \frac1k\sum_{n\ge1}G_\kappa(n/k)
              \longrightarrow\int_0^\infty G_\kappa(x)\dd x.
 \label{eq:riemann-limit}
\end{equation}
Indeed, for $0<\delta<1$, the integral over $(0,\delta)$
has absolute value at most $\delta$, and the portion of the
normalized mesh sum with $n/k<\delta$ is bounded by
$\delta+1/k$. On the remaining compact interval ordinary
Riemann-sum convergence applies. First let $k\to\infty$
and then let $\delta\downarrow0$. This argument does not
require a limit of $G_\kappa$ at zero.

For every real $k\ge1$ there is the exact identity
\begin{equation}
 \sum_{n\ge1}J_\kappa(k/n)
        =k\harm_{\lfloor k\rfloor}
                         +\sum_{n\ge1}G_\kappa(n/k).
 \label{eq:harmonic-separation}
\end{equation}
Since $\harm_{\lfloor k\rfloor}=\log k+\gamma+o(1)$,
\eqref{eq:riemann-limit} proves the expansion with coefficient
$\gamma+\int_0^\infty G_\kappa(x)\dd x$.
The substitution $z=1/x$ gives
\[
 \int_0^\infty G_\kappa(x)\dd x
       =\int_0^\infty
          \frac{J_\kappa(z)-z\ind_{[1,\infty)}(z)}{z^2}\dd z.
\]
Convergence also follows directly: the integrand is zero
below the positive band bottom and is bounded in absolute
value by $z^{-2}$ for $z\ge1$.
This proves \eqref{eq:ids-sum-law} and
\eqref{eq:coefficient-definition}. Combining it with
Proposition~\ref{prop:reduction} proves the last assertion.
\end{proof}

\subsection{Continuity, strict monotonicity, and limits}

\begin{proposition}[Range of the coefficient]
\label{prop:coefficient-range}
The function $C$ has all the properties in
\eqref{eq:coefficient-range}.
\end{proposition}

\begin{proof}
For a fixed phase, the finite-cell forms have a fixed domain
as $\kappa$ varies. Point evaluation is bounded on $H^1$,
so their eigenvalues vary continuously at every finite
$\kappa\ge0$. At a fixed $\kappa>0$ and $z>0$, only
finitely many phases have an eigenvalue exactly at $z^2$,
by Lemma~\ref{lem:phase}. The phase counts therefore converge
almost everywhere when $\kappa$ varies to that value, and
\eqref{eq:cell-bounds} dominates them. Hence $J_\kappa(z)$
is continuous in $\kappa>0$.

On a compact positive interval of couplings, the increasing
band bottom $r(\kappa)$ gives a uniform gap at zero.
The $z^{-2}$ bound at infinity is also uniform.
Dominated convergence in \eqref{eq:coefficient-definition}
proves continuity of $C$.

If $0<\kappa_1<\kappa_2$, form ordering gives
$J_{\kappa_1}(z)\ge J_{\kappa_2}(z)$ for all $z$.
Moreover, $r(\kappa_1)<r(\kappa_2)$. On the nonempty
interval between these band bottoms, the second integrated
count is zero and the first is positive by
\eqref{eq:ids-gap}. Therefore
\begin{equation}
 C(\kappa_1)-C(\kappa_2)
     =\int_0^\infty
          \frac{J_{\kappa_1}(z)-J_{\kappa_2}(z)}{z^2}\dd z>0.
 \label{eq:strict-monotonicity}
\end{equation}

As $\kappa\to\infty$, the increasing cell forms impose
$u(0)=0$ in the limit, and quasiperiodicity then imposes
$u(\pi)=0$. Their limiting form domain is the dense
Dirichlet domain. Each fixed indexed eigenvalue increases
to the corresponding Dirichlet eigenvalue $j^2$.
This also follows directly by min--max: Dirichlet test
spaces give the upper bound $j^2$; any sequence of normalized
eigenvectors with bounded energy is compact in $L^2$ and
has endpoint values tending to zero, so the limiting
variational spaces lie in the Dirichlet domain.
Thus, for noninteger $z$,
\begin{equation}
 J_\kappa(z)\longrightarrow\lfloor z\rfloor
                         \quad(\kappa\to\infty).
 \label{eq:ids-hard-limit}
\end{equation}
The uniform gap for $\kappa$ bounded away from zero and
the integrable tail bound justify dominated convergence
in the coefficient integral. Its limiting value is
\begin{align}
 \int_1^\infty\frac{\lfloor z\rfloor-z}{z^2}\dd z
 &=\sum_{j=1}^\infty
         \left[\frac1{j+1}-\log\frac{j+1}{j}\right]
   =\gamma-1.
 \label{eq:dirichlet-integral}
\end{align}
The last equality follows by taking partial sums,
$\harm_{L+1}-1-\log(L+1)$, and letting $L\to\infty$.
Consequently $\lim_{\kappa\to\infty}C(\kappa)=2\gamma-1$.

As $\kappa\downarrow0$, finite-cell form convergence
gives the free spectrum. At a fixed $z>0$, a free fiber
has eigenvalue $z^2$ at only finitely many phases, so
the same dominated-convergence argument and the free
normalization in Lemma~\ref{lem:cell-bounds} imply
\begin{equation}
 J_\kappa(z)\longrightarrow z
                           \quad(\kappa\downarrow0).
 \label{eq:ids-soft-limit}
\end{equation}
The portion of the coefficient integral over $z\ge1$
is at least $-1$, by \eqref{eq:cell-bounds}. On
$0<z<1$ the integrand is nonnegative. For any fixed
$0<\delta<1$, \eqref{eq:ids-soft-limit} gives
\[
 \lim_{\kappa\downarrow0}\int_\delta^1
                \frac{J_\kappa(z)}{z^2}\dd z
             =\int_\delta^1\frac{\mathrm dz}{z}
             =\log(1/\delta).
\]
It follows that
$\liminf_{\kappa\downarrow0}C(\kappa)
\ge\gamma-1+\log(1/\delta)$ for every such $\delta$.
Letting $\delta\downarrow0$ proves the infinite endpoint
limit. Strict monotonicity and continuity now give the
exact range in \eqref{eq:coefficient-range}.
\end{proof}
